\Chapter{Założenia organizacyjne wdrożenia systemu EMR}
\label{chapter:implementation}

\section{Wprowadzenie}

Wdrożenie systemu EMR w Szpitalu im. Bajgla stanowi nie tylko projekt technologiczny, ale również dużą zmianę organizacyjną obejmującą sposób codziennej pracy personelu medycznego i administracyjnego. Oprócz przygotowania samego systemu konieczne jest również zaplanowanie sposobu wdrożenia rozwiązania, migracji dokumentacji oraz organizacji pracy użytkowników w okresie przejściowym.

Celem przedstawionych założeń jest ograniczenie problemów organizacyjnych i technicznych związanych z cyfryzacją placówki przy jednoczesnym zachowaniu ciągłości pracy szpitala. Założenia te mają również zmniejszyć ryzyko błędów podczas wdrożenia oraz ułatwić adaptację personelu do nowego środowiska pracy.

\section{Etapowe wdrażanie systemu}

Ze względu na skalę działalności Szpitala im. Bajgla wdrożenie systemu powinno odbywać się etapowo. Rozwiązanie to pozwala ograniczyć ryzyko organizacyjne oraz uniknąć sytuacji, w której jednoczesna zmiana sposobu pracy całej placówki mogłaby doprowadzić do problemów z funkcjonowaniem oddziałów.

W pierwszym etapie system powinien zostać wdrożony na wybranych oddziałach o mniejszym stopniu złożoności organizacyjnej. Pozwoli to przetestować działanie podstawowych funkcjonalności oraz wprowadzić ewentualne poprawki przed rozszerzeniem systemu na pozostałe jednostki szpitala. Dopiero po zakończeniu fazy pilotażowej rozwiązanie mogłoby zostać wdrożone na oddziałach wymagających szybkiego przepływu informacji, takich jak SOR lub oddział intensywnej terapii.

Migracja dokumentacji papierowej również powinna odbywać się stopniowo. W pierwszej kolejności do systemu przenoszone byłyby aktualne dane pacjentów oraz dokumentacja wykorzystywana w bieżącej pracy oddziałów. Archiwalna dokumentacja mogłaby być cyfryzowana etapowo w późniejszym okresie, co pozwoliłoby ograniczyć koszty i obciążenie organizacyjne związane z jednoczesnym przetwarzaniem bardzo dużej ilości danych.

W okresie przejściowym przewidziano równoczesne funkcjonowanie części dokumentacji papierowej i cyfrowej. Takie rozwiązanie umożliwia stopniowe przyzwyczajanie personelu do nowego systemu oraz ogranicza ryzyko utraty dostępu do danych podczas początkowego etapu wdrożenia.

\section{Szkolenie i organizacja pracy personelu}

Wdrożenie nowego systemu wymaga odpowiedniego przygotowania personelu medycznego oraz administracyjnego do pracy w środowisku cyfrowym. Z tego powodu przed rozpoczęciem korzystania z systemu konieczne jest przeprowadzenie szkoleń obejmujących podstawowe funkcjonalności oraz procedury pracy z dokumentacją elektroniczną.

Szkolenia powinny być dostosowane do zakresu obowiązków poszczególnych grup użytkowników. Lekarze, pielęgniarki, laboranci oraz pracownicy administracyjni korzystają z różnych elementów systemu, dlatego zakres przekazywanych informacji powinien odpowiadać rzeczywistym potrzebom danej grupy.

W celu ograniczenia problemów organizacyjnych przewidziano również przygotowanie podstawowych procedur korzystania z systemu. Dotyczy to między innymi sposobu dodawania dokumentacji, aktualizowania danych pacjentów, przekazywania informacji pomiędzy oddziałami oraz postępowania w sytuacjach awaryjnych.

Istotnym elementem organizacji pracy jest również określenie odpowiedzialności za wprowadzanie i aktualizowanie danych. Każdy użytkownik powinien posiadać indywidualne konto umożliwiające identyfikację działań wykonywanych w systemie. Rozwiązanie to pozwala ograniczyć ryzyko błędów oraz zwiększa bezpieczeństwo dokumentacji medycznej.

\section{Założenia techniczne i bezpieczeństwo systemu}

Projektowany system zakłada wykorzystanie centralnej bazy danych umożliwiającej przechowywanie dokumentacji pacjentów w jednym spójnym środowisku. Dzięki temu możliwa będzie synchronizacja danych pomiędzy oddziałami oraz dostęp do aktualnych informacji w czasie rzeczywistym.

Istotnym założeniem projektu jest zapewnienie odpowiedniego poziomu bezpieczeństwa danych medycznych. System powinien wykorzystywać indywidualne konta użytkowników, podział poziomów dostępu oraz szyfrowanie danych i transmisji informacji. Dzięki temu możliwe będzie ograniczenie dostępu do dokumentacji wyłącznie do osób posiadających odpowiednie uprawnienia.

W projekcie przewidziano także tworzenie regularnych kopii zapasowych oraz procedur odzyskiwania danych po awarii systemu. Rozwiązanie to ma ograniczyć ryzyko utraty dokumentacji medycznej oraz umożliwić szybkie przywrócenie działania systemu w przypadku problemów technicznych.

Dodatkowo system powinien umożliwiać monitorowanie aktywności użytkowników oraz rejestrowanie zmian wykonywanych w dokumentacji medycznej. Funkcjonalność ta jest istotna zarówno z perspektywy bezpieczeństwa danych, jak i zgodności z wymaganiami dotyczącymi ochrony informacji medycznych obowiązującymi w Unii Europejskiej.

\section{Postępowanie w sytuacjach awaryjnych}

Ze względu na całodobowy charakter pracy szpitala konieczne jest przygotowanie procedur umożliwiających funkcjonowanie placówki również w przypadku czasowej niedostępności systemu. W sytuacji awarii personel powinien posiadać możliwość przejścia na uproszczony tryb pracy oraz czasowego wykorzystania dokumentacji zastępczej do momentu przywrócenia pełnej funkcjonalności systemu.

W projekcie przewidziano również stałe wsparcie techniczne po wdrożeniu systemu. W początkowym okresie działania rozwiązania szczególnie ważne będzie szybkie reagowanie na problemy zgłaszane przez użytkowników oraz możliwość bieżącego wprowadzania poprawek organizacyjnych i technicznych.